% --- Chapter: Error Handling ---
\section{Error Handling}

CRISP provides a modern error handling system with \texttt{try}/\texttt{catch}/\texttt{finally} blocks, catchable \texttt{throw} and \texttt{die} keywords, and proper propagation of control-flow signals through error boundaries.

\subsection{Try/Catch/Finally}

The \texttt{try} block executes code that may produce errors. If an error occurs, the optional \texttt{catch} block receives it. The \texttt{finally} block always executes regardless of errors:

\begin{tcolorbox}[colback=codebg, colframe=darkpurple, colbacktitle=darkpurple, title=\textbf{\textcolor{codegold}{Basic Try/Catch}}, fonttitle=\bfseries, sharp corners]
\begin{lstlisting}
# Basic try/catch
try {
    let x = 1 / 0;
} catch e {
    say "Error: ", e;
}
# Output: Error: Division by zero

# Catch with alias syntax (catch e as err)
try {
    die("something went wrong");
} catch e as err {
    say "Caught: ", err;
}
# Output: Caught: something went wrong
\end{lstlisting}
\end{tcolorbox}

\begin{tcolorbox}[colback=codebg, colframe=darkpurple, colbacktitle=darkpurple, title=\textbf{\textcolor{codegold}{Try/Catch/Finally}}, fonttitle=\bfseries, sharp corners]
\begin{lstlisting}
# Finally always executes
try {
    say "trying...";
    die("oops");
} catch e {
    say "error: ", e;
} finally {
    say "cleanup always runs";
}
# Output:
# trying...
# error: oops
# cleanup always runs

# Finally runs even without catch
try {
    say "work";
} finally {
    say "done";
}
# Output: work / done

# Nested error handling
try {
    try {
        throw "inner error";
    } catch e {
        say "inner: ", e;
    }
} catch e {
    say "outer: ", e;
}
# Output: inner: inner error
\end{lstlisting}
\end{tcolorbox}

\subsection{Throw and Die}

Both \texttt{throw} and \texttt{die} produce catchable \texttt{UserError} values. They are interchangeable:

\begin{tcolorbox}[colback=codebg, colframe=darkpurple, colbacktitle=darkpurple, title=\textbf{\textcolor{codegold}{Throw and Die}}, fonttitle=\bfseries, sharp corners]
\begin{lstlisting}
# throw — catchable user error
try {
    throw "invalid configuration";
} catch e {
    say "Problem: ", e;
}

# die — identical behavior, also catchable
try {
    die "fatal condition";
} catch e {
    say "Recovered from: ", e;
}

# Uncaught throw/die terminates the script
throw "unrecoverable";
# Error: unrecoverable
\end{lstlisting}
\end{tcolorbox}

\subsection{Control Flow Through Try/Catch}

\texttt{return}, \texttt{break}, and \texttt{continue} propagate through \texttt{try/catch/finally} correctly — the \texttt{finally} block always executes before the signal continues:

\begin{tcolorbox}[colback=codebg, colframe=darkpurple, colbacktitle=darkpurple, title=\textbf{\textcolor{codegold}{Control Flow with Finally}}, fonttitle=\bfseries, sharp corners]
\begin{lstlisting}
# break respects finally
while true {
    try {
        say "inside";
        break;
    } finally {
        say "cleanup!";
    }
}
# Output: inside / cleanup!

# continue respects finally
for i in 0..3 {
    try {
        if i == 1 { continue; }
        say "processing ", i;
    } finally {
        say "iter ", i;
    }
}
# Output: processing 0 / iter 0 / iter 1 / processing 2 / iter 2

# return from function through finally
fn demo() {
    try {
        return "early";
    } finally {
        say "still runs!";
    }
}
say demo();
# Output: still runs! / early
\end{lstlisting}
\end{tcolorbox}

\subsection{Assert}

The \texttt{assert} function verifies a condition and produces an error if it fails:

\begin{tcolorbox}[colback=codebg, colframe=darkpurple, colbacktitle=darkpurple, title=\textbf{\textcolor{codegold}{Assert}}, fonttitle=\bfseries, sharp corners]
\begin{lstlisting}
# Basic assertion
assert(2 + 2 == 4);  # passes silently

# Assertion with custom message
let x = -1;
assert(x > 0, "x must be positive");
# Error: x must be positive
\end{lstlisting}
\end{tcolorbox}

\begin{tcolorbox}[colback=lightlavender, colframe=softvioletdark, title=\textbf{\textcolor{darkpurple}{Design Note: Error Signals}}, fonttitle=\bfseries, sharp corners]
CRISP implements four internal sentinel errors (\texttt{RuntimeError::ReturnSignal}, 
\texttt{BreakSignal}, \texttt{ContinueSignal}, \texttt{ExitSignal}) that unwind 
through the interpreter's call stack. These are not caught by \texttt{catch} blocks — 
only real errors are. The \texttt{try/catch/finally} handler distinguishes between them:
\begin{itemize}
    \item \textbf{Real errors:} Captured by \texttt{catch}, then \texttt{finally} executes
    \item \textbf{Control-flow signals:} Bypass \texttt{catch}, \texttt{finally} still executes, then the signal is re-emitted
    \item \textbf{ExitSignal:} Propagates through all boundaries including
    \texttt{try/catch} and loops. \texttt{finally} blocks execute, then the
    process terminates with the given exit code. Not catchable by \texttt{catch}.
\end{itemize}
This ensures \texttt{break} from inside a \texttt{try} block inside a loop works as expected — the loop receives the \texttt{BreakSignal} after \texttt{finally} runs.
\end{tcolorbox}

\subsection{Exit}

The \texttt{exit()} function terminates the process with an optional exit code. It produces a special \texttt{ExitSignal} that is \textbf{not catchable} by \texttt{try/catch} blocks. However, \texttt{finally} blocks still execute before the process terminates:

\begin{tcolorbox}[colback=codebg, colframe=darkpurple, colbacktitle=darkpurple, title=\textbf{\textcolor{codegold}{Exit}}, fonttitle=\bfseries, sharp corners]
\begin{lstlisting}
# Exit with default code 0
exit();

# Exit with specific code
exit(1);

# finally runs before exit
try {
    exit(2);
} finally {
    say "cleanup runs!";
}
# Output: cleanup runs!
# Process exits with code 2

# _exit() skips finally — immediate hard exit
_exit(0);  # no cleanup, process dies immediately
\end{lstlisting}
\end{tcolorbox}

\begin{tcolorbox}[colback=lightlavender, colframe=softvioletdark, title=\textbf{\textcolor{darkpurple}{Design Note: ExitSignal vs. \_exit}}, fonttitle=\bfseries, sharp corners]
CRISP provides two exit mechanisms:

\begin{itemize}
    \item \texttt{exit(code?)} — Emits \texttt{ExitSignal}, which propagates through \texttt{finally} blocks before the process terminates. This ensures cleanup code runs. Not catchable by \texttt{catch} — it always terminates.
    \item \texttt{\_exit(code?)} — Calls Rust's \texttt{std::process::exit} directly. No cleanup, no \texttt{finally}, immediate termination. Use when you need a hard stop.
\end{itemize}

The \texttt{ExitSignal} uses the same sentinel-error mechanism as \texttt{return}/\texttt{break}/\texttt{continue} to unwind the call stack, ensuring \texttt{finally} blocks execute before the interpreter propagates the signal to the top level and calls \texttt{std::process::exit}.
\end{tcolorbox}
